\section{Scope, provenance, and conventions}
\label{sec:scope}

The starting point is the ProveIt repository at revision
\begin{center}
\texttt{c2cb285b64273e4cb2c8d7a03c68a4a99ac812d4}.
\end{center}
The source inventory under
\src{Analysis/Polylogarithms/docs/manuscript} and all nine ZIP reports
present under \src{docs/incoming} at that revision were retrieved.
Selected canonical manuscript sections were read alongside the reports.
The focused reading concerned the integration, harmonic, reflection,
Stieltjes contact, and Tornheim sections. This is a continuation of those
parts, not a claim to have audited every proof in the four-volume project.
The accompanying \src{SOURCE_AUDIT.md} identifies the relevant files and
the exact boundary between inherited results and the present extensions.

The latest incoming reports matter mathematically. They already prove
the cubic digamma--harmonic bridge, the independent cubic Tornheim
formula, the polynomial reflected Stieltjes closure, and several
all-depth spectral completion results. Repeating any of these as a new
theorem would misrepresent the state of the project. The questions
continued here are more specific:
\begin{enumerate}
\item The nonpolynomial reflected generator requested in Q8 of
\emph{Exact Resonant Identities} \cite{RepoExact}.
\item The necessary and sufficient centered cancellation criterion for
generalized harmonic polynomials requested in Q11 of the same report.
\item Simultaneous and hierarchical geometric collisions beyond the
two-factor case, asked for in the periodic contact reports
\cite{RepoPeriodic,RepoMixed}.
\item The quartic directional layer and its cancellation map, requested
in \emph{Independent Orders and Exact Reductions}, Section 7,
and the cubic harmonic report \cite{RepoIndependent,RepoCubic}.
\end{enumerate}

These were posed as research questions. We do not recast them as named
classical conjectures. The results below resolve the stated analytic and
formal-algebraic components; the arithmetic reduction of the retained
constants is a separate problem.

\subsection{Basic functions and the meaning of a derivative}

We use
\[
 \zeta(s,a)=\sum_{n=0}^{\infty}(n+a)^{-s},\qquad
 \Li_s(z)=\sum_{n=1}^{\infty}\frac{z^n}{n^s}
\]
in their initial convergence domains, followed by the explicitly stated
analytic continuation. The generalized Stieltjes functions are defined by
\begin{equation}\label{scope:stieltjes}
 \zeta(1+u,x)=\frac1u+
       \sum_{m=0}^{\infty}\frac{(-1)^m}{m!}\gamma_m(x)u^m.
\end{equation}
In particular $\gamma_0(x)=-\psi(x)$, where
$\psi=\Gamma'/\Gamma$. Write $\gamma=\gamma_0(1)$ and
$\gamma_m=\gamma_m(1)$. A prime on $\psi$ or $\gamma_m$ means an
argument derivative. A superscript $\zeta^{(j)}(s)$ means an order
derivative of the Riemann zeta function. Multivariable subscripts on
the Tornheim remainder mean partial derivatives, always evaluated at
the point specified. These distinctions are essential at a spectral pole.

The generalized harmonic numbers are
\[
 H_n^{(r)}=\sum_{k=1}^n k^{-r},\qquad H_n=H_n^{(1)}.
\]
All powers of positive real variables use the real logarithm. Thus, for
complex $\lambda$, the expression $(\sin(\pi x)/\pi)^{2\lambda}$
on $0<x<1$ is unambiguous. Meromorphic identities are continued from
a nonempty open convergence domain; a value at a crossing of polar
divisors is never assigned by informal substitution.

\subsection{Coordinate finite parts}

For an endpoint expansion in real powers and integral logarithmic powers,
the unit-coordinate finite part is the coefficient independent of both
$\epsilon$ and $\log\epsilon$ in the cutoff integral. On the unit interval,
\begin{equation}\label{scope:FP}
 \FP\int_0^1 f(x)\dd x
   =\CT_{\epsilon\downarrow0}\int_\epsilon^{1-\epsilon}f(x)\dd x.
\end{equation}
The coordinates are $x$ at zero and $1-x$ at one. Independent endpoint
cutoffs give the same additive constant. For example,
\[
 \FP\int_0^1x^{-1}\log^j x\dd x=0,\qquad
 \FP\int_0^1x^{-r}\dd x=\frac1{1-r}\quad(r\ne1).
\]
At periodic seams, the same convention is applied on each side in the
unit argument coordinate. Geometric collision constants are then taken
in a second parameter after this first finite part has been formed.
The two operations can differ from first merging the factors.

All coefficient extractions below are from a germ proved holomorphic
after its full prescribed subtraction. Numerical quadrature is used
only to test identities already supplied with proofs. Exact finite
linear algebra establishes a statement about its specified relation
space, not arithmetic independence of the associated periods.

\subsection{Relation to the literature}

Hurwitz Fourier transforms, Beta and Gamma identities, and their
applications to logarithmic integrals are classical
\cite{DLMFHurwitz,DLMFBeta,EspinosaMoll}. Harmonic Stieltjes constants
and log-sine connections were studied before this project
\cite{Kargin,Can}. The Mellin continuation and Euler boundary relations
for Tornheim zeta are likewise established tools \cite{BorweinDilcher}.
The contributions claimed here are the particular completed generators,
the exact all-polynomial criterion, the explicit cluster subtraction,
and the higher directional reconstruction proved below. The source
search was focused; it does not establish worldwide priority for every
displayed specialization.
